\subsection{From synthetic to real data}
Three differences between the synthetic and the real data shaped the second case study. The first is that the measurement errors are unknown. Profiling them out of the likelihood gives a ranking criterion that needs no error model, but it removes the adequacy test, and with it the safeguard against accepting a law that is merely the best of a poor set. Proposed laws on real data were therefore treated as candidates for further experiments, even where they crossed the BIC threshold. Replicate samples, or error models from the validation of the analytical methods, would restore the test. The profiled criterion also weights each species by its own fitted error, so a species that is fitted almost exactly, such as a substrate measured as zero after its depletion, carries much weight; a floor on the error standard deviations would limit this effect.

The second difference is model structure. The lumped glucose equivalent cannot represent the diauxic batch phase, in which the yeast first grows on glucose and then on the ethanol it produced, and every candidate compensated with a higher initial biomass. Because all candidates share this growth model, the error affects the ranking of the conversion mechanisms less than it affects the growth parameters, which should not be interpreted physiologically. The symbolic regression hinted at the problem, proposing growth laws that saturate more weakly than the library allows.

The third difference is run-to-run variability. Joint fitting resolved the ambiguity of single runs in case study~1 because all runs shared their kinetics by construction. The real runs agreed on the mechanism, but not on its parameters, and a joint fit with shared parameters described some runs much worse than their individual fits (Section~\ref{sec:c2:joint}). The power of joint fitting to discriminate between mechanisms therefore depends on the reproducibility of the cultures, and the shared-kinetics assumption should be tested before it is used, as was done here at no extra cost. Mixed-effects formulations, with run-specific parameters distributed around population values, are a natural way to combine runs that share a mechanism but not exact parameter values. Despite these differences, the analyses of the real runs converged on a consistent picture: every run, analysed alone, selected the mechanism with substrate and product inhibition of the conversion and cell death, and the laws proposed by the symbolic regression for every analysis contained death and described the slowing of the conversion through xylitol. Experiments that vary xylitol independently of time and xylose, for example by adding xylitol to the batch medium, would separate the proposed exponential product inhibition from the hyperbolic one of the library.
